% Compile twice: lualatex 20001329.tex
\documentclass[11pt]{article}
\usepackage[margin=25mm]{geometry}
\usepackage{fontspec,unicode-math}
\setmainfont{Latin Modern Roman}
\setmathfont{Latin Modern Math}
\usepackage{amsmath,amssymb,mathtools}
\usepackage{xurl,hyperref}
\hypersetup{colorlinks=true,urlcolor=blue,linkcolor=blue}
\setcounter{secnumdepth}{0}
\setlength{\parindent}{0pt}
\setlength{\parskip}{5pt}
\setlength{\emergencystretch}{3em}
\newcommand{\sref}[2]{\hyperlink{source:#1:#2}{[S#2]}}
\newcommand{\cataloguescope}[1]{\paragraph{Source question.}#1}
\newcommand{\floor}[1]{\left\lfloor#1\right\rfloor}
\begin{document}
% UnsolvedMath ID 20001329; source catalogue code AIM-CRYPTOGRAPHY-0016
\setcounter{section}{20001328}
\section{Multiple marked inputs: exclusion bounds and tight query regimes}\label{20001329}
{\small\sffamily UnsolvedMath ID 20001329 · Computer Science}
\subsection{Definitions and mathematical statement}

\paragraph{Shared notation.}Standard number systems are denoted by $\mathbb N,\mathbb Z,\mathbb Q,\mathbb R,\mathbb C$; index ranges, computational models and specialized notation follow the cited problem source and the definitions below.


An input is marked when $f(x)=1$. The requested output is m pairwise distinct marked inputs; repeating an earlier output does not count. Query complexity measures calls to f, with the domain size and number of marked inputs affecting the search cost.
\paragraph{Statement recorded in the supplied source.}
Find $x,y$ distinct with $f(x) = f(y) = 1$. Consider the case where there are many solutions at low density. How quickly can this be done? Faster than Grover?

More generally, find distinct $x_1,\dots, x_m$ such that $f(x_1) = \dots = f(x_m) = 1$. Can you do better than $m$ preimage searches? \sref{20001329}{2}
\paragraph{Scope qualification.}
The source emphasizes many low-density solutions but does not promise an exact marked-input count or specify whether it is known.
\paragraph{Source status.}
The supplied catalogue labels this record partially solved. Its source-status metadata was checked against the supplied snapshot on 2026-10-11. This import does not independently establish current open status or mathematical correctness.
\subsection{Short English statement}
Can several distinct marked inputs be found more efficiently than by carrying out separate quantum searches?
\subsection{Sources}
\begin{itemize}
\item[\hypertarget{source:20001329:1}{S1}] Ulam AI and the UnsolvedMath Contributors, UnsolvedMath: A Curated Collection of Open Mathematics Problems, record 20001329 (AIM-CRYPTOGRAPHY-0016). Catalogue curation and metadata: CC BY 4.0. The supplied source question is retained; definitions and the English summary are editorial additions.
\url{https://huggingface.co/datasets/ulamai/UnsolvedMath}.
\item[\hypertarget{source:20001329:2}{S2}] AIM Workshop Problem Lists, original problem source linked by catalogue record AIM-CRYPTOGRAPHY-0016
\url{http://aimpl.org/quantumalg/2/}.
\end{itemize}
\paragraph{Preparation and review.}
Original question attributed to the cited source and UnsolvedMath contributors. Codex (GPT-6; exact serving revision not exposed) assisted with definitions, English summaries, source inspection and formatting. No new human mathematical authorship or independent proof verification is asserted. Canonical ID 20001329 and source code AIM-CRYPTOGRAPHY-0016 checked against the complete supplied registry and existing repository statements; known flagged repetitions and obvious identical conjectures excluded. Full mathematical-equivalence/literature review has not been performed. Statement preparation and final review followed contribution prompts C and D; automated checks verify structure and identity, not open status.
\label{end:20001329}
\end{document}
